\section{Conceptual Foundation and Practical Extension}

\subsection{Understanding the bullwhip effect}

Our group considered it important to first examine the article by Lee, Padmanabhan, and Whang (1997), as it provides the conceptual basis needed to understand the article \textit{Machine Learning Demand Forecasting and Supply Chain Performance}. The authors define the \textit{bullwhip effect} as the amplification of order variability as one moves upstream in the supply chain. Their work is essential because it explains the structural origins of this phenomenon and highlights solutions such as better information sharing and improved coordination among supply chain actors.

\subsection{A practical response through forecasting}

In this context, Feizabadi's article can be seen as an operational extension of this theory. It shows how advanced forecasting methods based on machine learning can improve the quality of the demand signal transmitted across the supply chain. More specifically, the study suggests that hybrid models combining time series and explanatory variables can improve forecasting accuracy and enhance supply chain performance.

\subsection{Why both articles should be read together}

Understanding the bullwhip effect is therefore necessary to connect the theoretical source of supply chain disruption with the practical solution proposed in the more recent article. In other words, Lee, Padmanabhan, and Whang explain the problem, while Feizabadi provides a modern forecasting-based response to reduce its effects.
